
\documentclass{article}
\usepackage{colortbl}
\usepackage{makecell}
\usepackage{multirow}
\usepackage{supertabular}

\begin{document}

\newcounter{utterance}

\centering \large Interaction Transcript for game `eqbench', experiment `eq\_bench', episode 128 with reflexion{-}llama{-}3.1{-}8b{-}instruct.
\vspace{24pt}

{ \footnotesize  \setcounter{utterance}{1}
\setlength{\tabcolsep}{0pt}
\begin{supertabular}{c@{$\;$}|p{.15\linewidth}@{}p{.15\linewidth}p{.15\linewidth}p{.15\linewidth}p{.15\linewidth}p{.15\linewidth}}
    \# & \multicolumn{2}{c}{Player} && \multicolumn{2}{c}{Game Master} \\
    \hline

    \theutterance \stepcounter{utterance}  
    & & & \multicolumn{4}{p{0.6\linewidth}}{
        \cellcolor[rgb]{0.9,0.9,0.9}{
            \makecell[{{p{\linewidth}}}]{
                \texttt{\tiny{[P1$\langle$GM]}}
                \texttt{You have attempted this same task before and did not succeed. These are the plans you wrote after those attempts.} \\
\texttt{Trial \#1: New plan: I overestimated Olivia's acceptance of Dr. Fleischer's approach, as she explicitly stated she wasn't there for a philosophical discussion. I should have considered her annoyance and resentment more strongly, as she seems to be frustrated with the conversation's direction. Next time, I will give more weight to Olivia's direct statements and consider her emotional responses more nuanced, taking into account her resistance to Dr. Fleischer's approach. I will also consider the context of the conversation, where Olivia is likely feeling overwhelmed or stuck, and Dr. Fleischer's approach may be perceived as insensitive or dismissive.} \\
\\ 
\texttt{Trial \#2: New plan: I overestimated Olivia's acceptance of Dr. Fleischer's approach, as she explicitly stated she wasn't there for a philosophical discussion. I should have considered her annoyance and resentment more strongly, as she seems to be frustrated with the conversation's direction. However, I also underestimated her level of frustration, as I gave Indifference a score of 2. Next time, I will give more weight to Olivia's direct statements and consider her emotional responses more nuanced, taking into account her resistance to Dr. Fleischer's approach. I will also consider the context of the conversation, where Olivia is likely feeling overwhelmed or stuck, and Dr. Fleischer's approach may be perceived as insensitive or dismissive. To improve my scoring, I will pay closer attention to Olivia's language and tone, and consider the emotional undertones of her statements. I will also consider the power dynamics at play in the conversation, where Dr. Fleischer is in a position of authority and Olivia may feel pressured to conform to his expectations.} \\
\\ 
\\ 
\\ 
\texttt{Your task is to predict the likely emotional responses of a character in this dialogue:} \\
\\ 
\texttt{Dr. Fleischer: You've been doing well, Olivia. Your writing's been improving, hasn't it?} \\
\texttt{Olivia: I suppose. But that's not why I'm here, is it?} \\
\texttt{Dr. Fleischer: No, it's not. But it's all connected, isn't it? Your writing, your past, your anxieties...} \\
\texttt{Olivia: I'm not here for a philosophical discussion, Dr. Fleischer.} \\
\texttt{Dr. Fleischer: Of course, you're not. But sometimes, understanding the connections helps us untangle the knots, wouldn't you agree?} \\
\texttt{Olivia: Or it just makes things more complicated.} \\
\texttt{Dr. Fleischer: Perhaps. But isn't it better to face the complexity than to ignore it?} \\
\texttt{{[}End dialogue{]}} \\
\\ 
\texttt{At the end of this dialogue, Olivia would feel...} \\
\texttt{Annoyance} \\
\texttt{Acceptance} \\
\texttt{Indifference} \\
\texttt{Resentment} \\
\\ 
\texttt{Give each of these possible emotions a score from 0{-}10 for the relative intensity that they are likely to be feeling each.} \\
\\ 
\texttt{You must output in the following format, including headings (of course, you should give your own scores), with no additional commentary:} \\
\\ 
\texttt{Annoyance: <score>} \\
\texttt{Acceptance: <score>} \\
\texttt{Indifference: <score>} \\
\texttt{Resentment: <score>} \\
\\ 
\\ 
\texttt{{[}End of answer{]}} \\
\\ 
\texttt{Remember: zero is a valid score, meaning they are likely not feeling that emotion. You must score at least one emotion > 0.} \\
\\ 
\texttt{Your answer:} \\
            }
        }
    }
     \\ \\

    \theutterance \stepcounter{utterance}  
    & \multicolumn{4}{p{0.6\linewidth}}{
        \cellcolor[rgb]{0.561,0.737,0.561}{
            \makecell[{{p{\linewidth}}}]{
                \texttt{\tiny{[P1$\rangle$GM]}}
                \texttt{Annoyance: 8} \\
\texttt{Acceptance: 0} \\
\texttt{Indifference: 2} \\
\texttt{Resentment: 6} \\
            }
        }
    }
    & & \\ \\

    \theutterance \stepcounter{utterance}  
    & & & \multicolumn{2}{p{0.3\linewidth}}{
        \cellcolor[rgb]{0.9,0.9,0.9}{
            \makecell[{{p{\linewidth}}}]{
                \texttt{\tiny{[GM$|$GM]}}
                \texttt{Annoyance: 8} \\
\texttt{Acceptance: 0} \\
\texttt{Indifference: 2} \\
\texttt{Resentment: 6} \\
            }
        }
    }
    & & \\ \\

    \theutterance \stepcounter{utterance}  
    & & & \multicolumn{2}{p{0.3\linewidth}}{
        \cellcolor[rgb]{0.9,0.9,0.9}{
            \makecell[{{p{\linewidth}}}]{
                \texttt{\tiny{[GM$|$GM]}}
                \texttt{Annoyance: 7} \\
\texttt{Acceptance: 0} \\
\texttt{Indifference: 0} \\
\texttt{Resentment: 3} \\
            }
        }
    }
    & & \\ \\

    \theutterance \stepcounter{utterance}  
    & & & \multicolumn{2}{p{0.3\linewidth}}{
        \cellcolor[rgb]{0.9,0.9,0.9}{
            \makecell[{{p{\linewidth}}}]{
                \texttt{\tiny{[GM$|$GM]}}
                \texttt{game\_result = LOSE} \\
            }
        }
    }
    & & \\ \\

    \theutterance \stepcounter{utterance}  
    & & \multicolumn{4}{p{0.6\linewidth}}{
        \cellcolor[rgb]{0.95,0.95,0.95}{
            \makecell[{{p{\linewidth}}}]{
                \texttt{\tiny{[P1$\leftrightarrow$P]}}
                \texttt{Post{-}episode Reflexion (not sent to GM):} \\
\texttt{New plan: I overestimated Olivia's acceptance of Dr. Fleischer's approach, as she explicitly stated she wasn't there for a philosophical discussion. I should have considered her annoyance and resentment more strongly, as she seems to be frustrated with the conversation's direction. However, I also underestimated her level of frustration, as I gave Indifference a score of 2. Next time, I will give more weight to Olivia's direct statements and consider her emotional responses more nuanced, taking into account her resistance to Dr. Fleischer's approach. I will also consider the context of the conversation, where Olivia is likely feeling overwhelmed or stuck, and Dr. Fleischer's approach may be perceived as insensitive or dismissive. To improve my scoring, I will pay closer attention to Olivia's language and tone, and consider the emotional undertones of her statements. I will also consider the power dynamics at play in the conversation, where Dr. Fleischer is in a position of authority and Olivia may feel pressured to conform to his expectations. Additionally, I will analyze the conversation more closely to identify specific moments where Olivia's emotions are triggered, such as when Dr. Fleischer says "But it's all connected, isn't it?" which may be perceived as dismissive or condescending. By paying closer attention to these moments, I can develop a more accurate understanding of Olivia's emotional state and provide a more nuanced scoring.} \\
            }
        }
    }
    & \\ \\

\end{supertabular}
}

\end{document}
